\section{Separable Triplet Reduction and Probabilistic Error Propagation}
\label{sec:analysis}

The tensor structure reduces each isolated multidimensional triplet exactly to the scalar Candan problem. Interference and noise then enter as explicit perturbations controlled by Fourier coherence, dynamic range, and SNR.

\textbf{Proposition 1 (isolated separable triplet reduction):}
For component $\ell$, let $b_{r,\ell}=k_{r,\ell}+\delta_{r,\ell}$ with $\delta_{r,\ell}\in[-1/2,1/2]$. Without other components or noise, for axis $d$ and $q\in\{-1,0,1\}$,
\begin{equation}
\begin{aligned}
 Z^{(\ell)}[\mathbf k_\ell+q\mathbf e_d]
 &=C_{-d,\ell}g_{N_d}(\delta_{d,\ell}-q),\\
 C_{-d,\ell}&=\beta_\ell\prod_{r\ne d}g_{N_r}(\delta_{r,\ell}).
\end{aligned}
 \label{eq:isolated_triplet_factorization}
\end{equation}
The common complex factor cancels from the quotient, giving
\begin{equation}
 \frac{Z_{d,-}^{(\ell)}-Z_{d,+}^{(\ell)}}
 {2Z_{d,0}^{(\ell)}-Z_{d,-}^{(\ell)}-Z_{d,+}^{(\ell)}}
 =\frac{\tan(\pi\delta_{d,\ell}/N_d)}{\tan(\pi/N_d)}.
 \label{eq:isolated_candan_ratio}
\end{equation}
Hence the unclipped update is
\begin{equation}
 \phi_{N_d}(\delta_{d,\ell})
 =\frac{\tan(\pi\delta_{d,\ell}/N_d)}{\pi/N_d}.
 \label{eq:candan_finite_map}
\end{equation}
\emph{Proof:} Equation~\eqref{eq:tensor_fft_factor} gives the product in \eqref{eq:isolated_triplet_factorization}; changing only axis $d$ leaves all other factors unchanged. Substitution of \eqref{eq:complex_dirichlet} and the angle-addition identities yields \eqref{eq:isolated_candan_ratio}. \hfill$\square$

Thus separability removes every non-updated axis from the isolated quotient. The finite-length map retains the small bias
\begin{equation}
 b_N(\delta)=|\phi_N(\delta)-\delta|=O(N^{-2}).
 \label{eq:candan_finite_bias}
\end{equation}

\textbf{Proposition 2 (multi-component Gaussian quotient bound):}
Assume $N_d\geq3$, independent tensor noise $\mathcal{CN}(0,\sigma^2)$, and a fixed selected-bin/component association. Write
\begin{equation}
 Z_{d,q}=C_{-d,\ell}g_{N_d}(\delta_{d,\ell}-q)+u_{d,q}+w_{d,q},
 \label{eq:triplet_signal_interference_noise}
\end{equation}
where
\begin{equation}
 u_{d,q}=\sum_{j\ne\ell}\beta_j\prod_{r=1}^{D}g_{N_r}
 (b_{r,j}-k_{r,\ell}-q\mathbf 1_{\{r=d\}}).
\end{equation}
Define
\begin{equation}
 I_{d,q,\ell}=\sum_{j\ne\ell}|\beta_j|\prod_{r=1}^{D}\mu_{N_r}
 (b_{r,j}-k_{r,\ell}-q\mathbf 1_{\{r=d\}}),
 \label{eq:interference_envelope}
\end{equation}
and the target numerator, denominator, and quotient
\begin{align}
 n_{d,\ell}&=C_{-d,\ell}\{g_{N_d}(\delta+1)-g_{N_d}(\delta-1)\},\nonumber\\
 d_{d,\ell}&=C_{-d,\ell}\{2g_{N_d}(\delta)-g_{N_d}(\delta+1)-g_{N_d}(\delta-1)\},\nonumber\\
 \gamma_{d,\ell}&=n_{d,\ell}/d_{d,\ell}
 =\tan(\pi\delta/N_d)/\tan(\pi/N_d),
 \label{eq:target_num_den}
\end{align}
where $\delta=\delta_{d,\ell}$. For $\eta\in(0,1)$ let
\begin{align}
 B_{n,d,\ell}&=I_{d,-1,\ell}+I_{d,+1,\ell}
 +\sigma\sqrt{2\log(2/\eta)},\nonumber\\
 B_{d,d,\ell}&=2I_{d,0,\ell}+I_{d,-1,\ell}+I_{d,+1,\ell}
 +\sigma\sqrt{6\log(2/\eta)}.
 \label{eq:explicit_bn_bd}
\end{align}
If $B_{d,d,\ell}<|d_{d,\ell}|$, then with probability at least $1-\eta$,
\begin{equation}
\begin{aligned}
 |\widehat\delta_{d,\ell}-\delta_{d,\ell}|
 &\le b_{N_d}(\delta_{d,\ell})\\
 &\quad+c_{N_d}\frac{B_{n,d,\ell}+|\gamma_{d,\ell}|B_{d,d,\ell}}
 {|d_{d,\ell}|-B_{d,d,\ell}},\\[-1mm]
 c_N&=\frac{\tan(\pi/N)}{\pi/N}.
\end{aligned}
 \label{eq:coordinate_error_bound}
\end{equation}
\emph{Proof:} The unitary FFT preserves the noise law. Numerator noise is $\mathcal{CN}(0,2\sigma^2)$ and denominator noise is $\mathcal{CN}(0,6\sigma^2)$. Complex-Gaussian tails and a union bound give \eqref{eq:explicit_bn_bd}. The quotient identity
\begin{equation*}
 \left|\frac{n+e_n}{d+e_d}-\frac nd\right|
 \le\frac{|e_n|+|n/d||e_d|}{|d|-|e_d|}
\end{equation*}
followed by the finite-length factor and nonexpansive clipping proves \eqref{eq:coordinate_error_bound}. \hfill$\square$

The dependence is explicit. If $\Gamma_\ell=\max_{j\ne\ell}|\beta_j|/|\beta_\ell|$ and $\bar\mu_{d,q,\ell}$ is the largest coherence product in \eqref{eq:interference_envelope}, then
\begin{equation}
 I_{d,q,\ell}\le |\beta_\ell|(L-1)\Gamma_\ell\bar\mu_{d,q,\ell}.
 \label{eq:dynamic_range_interference}
\end{equation}
Under the executed SNR convention, $\sigma=\sqrt{P_{\mathcal Y_0}}10^{-\mathrm{SNR}/20}$. A simultaneous statement for all $DL$ triplets replaces $\eta$ by $\eta/(DL)$. The observable statistic $\zeta$ in \eqref{eq:candan_stability} is the measured counterpart of the denominator margin; empirically, its Spearman association with clipping is $-0.525$, $-0.508$, and $-0.506$ on CDL-A/C/D.

\textbf{Corollary 1 (propagation to the joint fit):}
Let $\mathbf A_0$ and $\widehat{\mathbf A}$ contain normalized atoms at the true and refined coordinates. If $|\widehat b_{d,\ell}-b_{d,\ell}|\le\epsilon_{d,\ell}$ periodically, define
\begin{equation}
\begin{aligned}
 \Lambda_N&=\frac{2\pi}{N}\sqrt{\frac{(N-1)(2N-1)}{6}},\\
 E_A&=\left[\sum_{\ell=1}^{L}\left(\sum_{d=1}^{D}
 \Lambda_{N_d}\epsilon_{d,\ell}\right)^2\right]^{1/2}.
\end{aligned}
 \label{eq:dictionary_lipschitz}
\end{equation}
Then $\|\mathbf A_0-\widehat{\mathbf A}\|_2\le E_A$. If $\widehat{\mathbf A}$ has full column rank,
\begin{align}
 \|\widehat{\boldsymbol\beta}-\boldsymbol\beta\|_2
 &\le\frac{E_A\|\boldsymbol\beta\|_2+\|\mathbf w\|_2}
 {\sigma_{\min}(\widehat{\mathbf A})},\nonumber\\
 \|\widehat{\mathbf y}-\mathbf y_0\|_2
 &\le E_A\|\boldsymbol\beta\|_2+\|\mathbf w\|_2.
 \label{eq:coordinate_to_ls_bound}
\end{align}
\emph{Proof:} The derivative norm of \eqref{eq:steering_atom} is $\Lambda_N$. A path integral bounds each separable atom perturbation by $\sum_d\Lambda_{N_d}\epsilon_{d,\ell}$; the Frobenius bound gives $E_A$, followed by the standard joint-LS perturbation identities. \hfill$\square$

Finally, \eqref{eq:observable_residual} is exactly the received-energy difference between two fitted subspaces by the Pythagorean identity. This identity connects the triplet-level refinement to the scene-level return rule.
